\section{L\'imite local y clasificaci\'on modular del sistema infinito}

En una UHF infinita no existe, en general, una matriz densidad traza-clase
global dentro del \`algebra. El objeto correcto es la red
\begin{equation}
\label{rm:eq:modular-local-net}
\HHM^{(m)}=-\log\rho_m
\end{equation}
con restricciones de estados compatibles, o el generador modular en la
representaci\'on GNS. La expectativa tracial
\(\mathrm{id}\otimes\tau_9\) no preserva autom\'aticamente un Gibbs no
tracial y no sustituye la compatibilidad de estados.

\begin{theorem}[Tipo modular bajo persistencia fiel de ambos canales]
\label{rm:thm:modular-typeIII}
Sup\'ongase que:
\begin{enumerate}[label=(\roman*)]
\item cada corte es un producto de estados qutrit fieles con pesos
\(w_{90}\) y \(w_{120}\);
\item las inclusiones preservan el estado y la etiqueta de canal;
\item aparecen infinitamente muchos factores de cada tipo, ambos con
densidad asint\'otica positiva;
\item no se introducen cocientes espectrales adicionales y la
representaci\'on producto es factorial.
\end{enumerate}
Entonces el grupo aditivo de logaritmos de cocientes de autovalores es
\begin{equation}
\label{rm:eq:modular-ratio-group}
90\Dmod\Z+120\Dmod\Z=30\Dmod\Z.
\end{equation}
El cierre GNS ITPFI es el factor hiperinfinito
\begin{equation}
\label{rm:eq:modular-typeIII-lambda}
\boxed{
III_{\lambda_\partial},
\qquad
\lambda_\partial=\ee^{-30\Dmod}
=0.9966696464689573786108299769\ldots .}
\end{equation}
\end{theorem}

\begin{proof}
En cada enlace los cocientes de autovalores son potencias de
\(\ee^{-w_m}\). La persistencia infinita de ambos tipos genera el subgrupo
\(w_{90}\Z+w_{120}\Z\). Como \(\gcd(90,120)=30\), se obtiene
\cref{rm:eq:modular-ratio-group}. El criterio de raz\'on asint\'otica de
Araki--Woods para el producto fiel da el factor ITPFI hiperinfinito con
par\'ametro \(\ee^{-30\Dmod}\)
\cite{rm:ArakiWoods1968,rm:Takesaki}.
\end{proof}

El teorema declara sus hip\'otesis reales. Si desaparece uno de los canales,
si el producto no es compatible o si cambia el grupo de razones, debe
recalcularse el tipo.
